\section{Introduction}

Portfolio management involves allocating funds among different categories of assets to ensure the proper achievement of a balance between expected return and risk. Modern portfolio theory provides the base theoretical foundation for this process, arguing that investors should classify and evaluate securities portfolios not solely on their individuality, but also on expected returns and volatility, as well as their covariance with other assets in the portfolio, as argued by \textcite{markowitz1952}. Thus, diversification may assist investors in reducing portfolio risk without an equivalent reduction in expected return when selected assets are imperfectly correlated.

This case study sheds light on evaluating a diversified portfolio consisting of three Philippine equities: PLDT Inc. (\texttt{TEL}), Manila Electric Company (\texttt{MER}), and Aboitiz Equity Ventures (\texttt{AEV}); two international technology equities: NVIDIA Corporation (\texttt{NVDA}) and Meta Platforms, Inc. (\texttt{META}); and one alternative digital asset: Bitcoin (\texttt{BTC}). Based on these assets, the portfolio combines domestic Philippine securities, Western growth-oriented technology stocks, and a cryptocurrency. This allocation of diverse asset classes provides a structured opportunity to assess whether diversification across different platforms, industries, asset classes, and geographies enhances portfolio performance. However, each individual asset must be assessed carefully as certain assets exhibit substantial volatility and diversification benefits vary across market regimes and periods of financial stress \parencite{almeida2023}.

Through the use of daily adjusted closing price data from January 2, 2020 through December 31, 2025, this study gauges the return and risk profile of the investment portfolio benchmarked against the SPDR S\&P 500 ETF Trust (\texttt{SPY}), which reflects the performance of large-capitalization equities from the United States. Furthermore, daily log return measurements are calculated to examine descriptive statistics, return distributions, correlations, and volatility among the selected assets. The empirical analysis compares portfolio execution with and without monthly rebalancing to assess whether restoring portfolio weights periodically affects portfolio growth, risk exposure, and overall performance.

Moreover, the portfolio is evaluated using the Philippine 91-Day Treasury Bill as the proxy for the risk-free rate. This rate supports the computation of the Sharpe ratio, measuring excess return earned per unit of total risk \parencite{sharpe1966}. In addition, mean-variance optimization is applied to construct the efficient frontier and identify both the Global Minimum Variance portfolio and the Maximum Sharpe Ratio portfolio. These optimization outputs are compared with alternative asset allocations to determine which strategy offers the most favorable risk-return trade-off across the sample period.

Additionally, because standard deviation and average returns cannot fully capture extreme losses, this study conducts a comprehensive downside risk assessment. Non-parametric Historical Value at Risk, Expected Shortfall (Conditional Value at Risk), Parametric Gaussian Value at Risk, Monte Carlo simulations, and crisis stress testing are utilized to evaluate potential portfolio losses under normal and adverse market conditions. Expected Shortfall is particularly essential because it measures the conditional average loss beyond the Value at Risk threshold, providing deeper insight into severe tail-risk exposure \parencite{acerbi2002}.

Ultimately, the study aims to develop an evidence-based investment recommendation for an institutional investor considering exposure to Philippine equities, Bitcoin, and international technology equities. The recommendations synthesize empirical return distributions, correlation structures, optimized frontier allocations, benchmark-relative performance, and portfolio resilience across historical stress regimes, establishing a defensible balance between return potential, total volatility, and downside risk tolerance.

\subsection{Mandate and Assigned Asset Universe}

This report delivers a quantitative portfolio management consultation for Group 3 (Charlie). The assigned asset universe comprises six multi-asset constituents across four distinct asset classes:

\begin{enumerate}
  \item \textbf{Philippine Equities (Domestic Core):} PLDT Inc. (\texttt{TEL}), Manila Electric Company (\texttt{MER}), and Aboitiz Equity Ventures (\texttt{AEV}).
  \item \textbf{International Equities (Growth Satellite):} NVIDIA Corporation (\texttt{NVDA}) and Meta Platforms, Inc. (\texttt{META}).
  \item \textbf{Digital Assets (Alternative Satellite):} Bitcoin (\texttt{BTC}).
  \item \textbf{Benchmark ETF:} SPDR S\&P 500 ETF Trust (\texttt{SPY}).
  \item \textbf{Risk-Free Rate Baseline:} Philippine 91-Day Treasury Bill secondary market rate ($R_f = 5.25\%$ per annum benchmark).
\end{enumerate}

\subsection{Risk-Free Rate Benchmark Specification}

The hurdle rate for all excess return, Sharpe ratio, and tangency portfolio calculations is anchored to the Philippine 91-Day Treasury Bill (3-Month T-Bill) secondary market yield. Quoted at an annualized benchmark rate of $R_f = 5.25\%$, it is converted into a daily log-equivalent rate assuming a 252-trading-day annual convention:

\begin{equation}
  r_{f,\text{daily}} = \frac{\ln(1 + R_f)}{252} = \frac{\ln(1 + 0.0525)}{252} \approx 0.00020307 \quad (0.0203\% \text{ per day})
\end{equation}

This hurdle rate represents the riskless baseline return earned by domestic institutional capital.

\subsection{Report Organization}

The remainder of this report is structured as follows: Section 2 details data acquisition, FX conversion, calendar alignment, and missing value treatment. Section 3 conducts exploratory financial analysis and distributional diagnostics. Section 4 evaluates Buy-and-Hold versus Monthly Rebalanced equal-weighted portfolios. Section 5 solves for the Efficient Frontier, Global Minimum Variance (GMV), and Maximum Sharpe Ratio portfolios. Section 6 assesses downside risk via VaR, Expected Shortfall, Monte Carlo simulations, and stress testing. Section 7 presents executive investment recommendations, and Section 8 lists references.
